\documentclass[12pt,letterpaper, onecolumn]{exam}
\usepackage{amsmath}
\usepackage{amssymb}
\usepackage[lmargin=71pt, tmargin=1.2in]{geometry}
\usepackage{xcolor}
\usepackage[brazil]{babel}
\usepackage{natbib}
\usepackage{enumitem}
\usepackage{booktabs}
\usepackage[hidelinks]{hyperref}
\urlstyle{same}

\renewenvironment{solution}
  {\par\noindent\textbf{R:}\enspace\color{blue}}
  {\par\color{black}}

\thispagestyle{empty}

\begin{document}

\begingroup
    \centering
    \LARGE Lista 3 Gabarito - Econometria I - 2026\\[0.5em]
    \large Prof: André Luiz Pereira Mancha \par
    \large Monitor: Tuffy Licciardi Issa \par
    \small \href{mailto:tuffy@usp.br}{tuffy@usp.br} \quad
    \href{https://github.com/Tuffyli}{github.com/tuffyli} \par
\endgroup
\rule{\textwidth}{0.4pt}
\pointsdroppedatright
\printanswers
\renewcommand{\solutiontitle}{\noindent\textbf{Ans:}\enspace}

\begin{questions}

   \question (3.5 \citep{wooldridge2016}) Em um estudo que relaciona a nota média em curso superior (\textit{GPA}) ao tempo gasto em várias atividades, você distribui uma pesquisa para vários estudantes. Os estudantes devem responder quantas horas eles despendem, em cada semana, em quatro atividades: estudo (\textit{study}), sono (\textit{sleep}), trabalho (\textit{work}) e lazer (\textit{leisure}). Toda atividade é colocada em uma das quatro categorias, de modo que, para cada estudante, a soma das horas nas quatro atividades deve ser igual a 168.
    \begin{enumerate}[label=(\alph*)]
    \item No modelo:\begin{equation*}
        colGPA = \beta_0 + \beta_1study +\beta_2sleep + \beta_3work + \beta_4 leisure + u,
    \end{equation*} faz sentido manter \textit{sleep}, \textit{work} e \textit{leisure} fixos, enquanto \textit{study} varia?
    \item Explique a razão desse modelo violar a Hipótese RLM.3.
    \item Como você poderia reformular o modelo de modo que seus parâmetros tivessem uma interpretação útil e ele satisfizesse a Hipótese RLM.3?
    \end{enumerate}

\begin{solution}
\begin{enumerate}[label=(\alph*)]
    \item Não. Como a soma das quatro atividades é sempre 168, se \textit{study} varia, pelo menos uma das outras categorias precisa variar também.
    \item O modelo viola a hipótese de ausência de colinearidade perfeita, pois
    \[
    study+sleep+work+leisure=168
    \]
    implica uma relação linear exata entre os regressors.
    \item Uma forma é omitir uma categoria, por exemplo \textit{leisure}, e estimar
    \[
    colGPA=\beta_0+\beta_1study+\beta_2sleep+\beta_3work+u.
    \]
\end{enumerate}
\end{solution}

\question (3.3 \citep{wooldridge2016}) Considere o modelo
\[
sleep=\beta_0+\beta_1totwrk+\beta_2educ+\beta_3age+u,
\]
em que \textit{sleep} e \textit{totwrk} são medidos em minutos por semana e \textit{educ} e \textit{age} em anos.
\begin{enumerate}[label=(\alph*)]
    \item Se os adultos escolhem dormir e trabalhar, qual é o sinal de $\beta_1$?
    \item Que sinais você espera que $\beta_2$ e $\beta_3$ terão?
    \item Usando os dados do arquivo \textit{SLEEP75}, a equação estimada é
    \[
    \widehat{sleep}=3638{,}25-0{,}148\,totwrk-11{,}13\,educ+2{,}20\,age
    \]
    com \(n=706\) e \(R^2=0{,}113\). Se alguém trabalha cinco horas a mais por semana, qual é a queda, em minutos, no valor esperado de \textit{sleep}? Esse valor representa uma mudança grande?
    \item Discuta o sinal e a magnitude do coeficiente de \textit{educ}.
    \item Você diria que \textit{totwrk}, \textit{educ} e \textit{age} explicam muito a variação de \textit{sleep}? Quais outros fatores poderiam afetar o tempo gasto dormindo? É provável que sejam correlacionados com \textit{totwrk}?
\end{enumerate}

\begin{solution}
\begin{enumerate}[label=(\alph*)]
    \item Espera-se \(\beta_1<0\): mais tempo de trabalho tende a reduzir o tempo disponível para dormir.
    \item Espera-se \(\beta_2<0\) e \(\beta_3>0\).
    \item Cinco horas a mais por semana equivalem a \(5\times 60=300\) minutos. Queda prevista é
    \[
    -0{,}148(300)=-44{,}4
    \]
    minutos por semana. Não parece uma mudança muito grande.
    \item O coeficiente de \textit{educ} é negativo: um ano adicional de educação está associado a \(11{,}13\) minutos a menos de sono por semana. A magnitude é pequena.
    \item Não muito. O \(R^2=0{,}113\) indica que apenas cerca de 11,3\% da variação de \textit{sleep} é explicada por essas variáveis. Outros fatores: estado de saúde, presença de filhos pequenos, estresse, tipo de ocupação, renda, problemas de sono. É possível que alguns deles sejam correlacionados com \textit{totwrk}.
\end{enumerate}
\end{solution}

\question (3.2 \citep{wooldridge2010}) Com os dados do arquivo \textit{WAGE2}, foi estimada a equação
\[
\widehat{educ}=10{,}36-0{,}094\,sibs+0{,}131\,meduc+0{,}210\,feduc
\]
com \(n=722\) e \(R^2=0{,}214\), em que \textit{educ} é anos de escolaridade, \textit{sibs} é o número de irmãos, \textit{meduc} é a escolaridade da mãe e \textit{feduc} é a escolaridade do pai.
\begin{enumerate}[label=(\alph*)]
    \item \textit{sibs} tem efeito esperado? Explique. Mantendo \textit{meduc} e \textit{feduc} fixos, quanto deveria \textit{sibs} aumentar para reduzir os anos previstos de educação formal em um ano?
    \item Discuta a interpretação do coeficiente de \textit{meduc}.
    \item Suponha que o Homem A não tenha irmãos e que sua mãe e seu pai tenham, cada um, 12 anos de educação formal. Suponha também que o Homem B não tenha irmãos e sua mãe e seu pai tenham, cada um, 16 anos de educação formal. Qual é a diferença prevista em anos de educação formal entre A e B?
\end{enumerate}

\begin{solution}
\begin{enumerate}[label=(\alph*)]
    \item Sim. Mais irmãos podem reduzir recursos familiares e reduzir a escolaridade esperada. Um aumento aproximado de
    \[
    \frac{1}{0{,}094}\approx 10{,}64
    \]
    irmãos.
    \item Um ano adicional de escolaridade da mãe aumenta a escolaridade prevista do filho em \(0{,}131\) ano.
    \item A diferença prevista é
    \[
    4(0{,}131)+4(0{,}210)=0{,}524+0{,}840=1{,}364.
    \]
    Portanto, B tem \(1{,}364\) anos de escolaridade prevista a mais do que A.
\end{enumerate}
\end{solution}

\question (3.3 \citep{wooldridge2010}) O salário incial (mediano) para recém-formados em Direito é determinado pela equação:\begin{equation*}
    \log(salary) = \beta_0 + \beta_1LSAT + \beta_2GPA + \beta_3\log(libvol) + \beta_4\log(cost) + \beta_5rank + u
\end{equation*} em que \textit{LSAT} é a nota mediana do LSAT (nota de ingresso no curso de Direito) dos recém-formados, \textit{GPA} é a nota mediana dos recém-formados nas disciplinas do curso de Direito, \textit{libvol} é o número de volumes da biblioteca da escola de Direito, \textit{cost} é o custo anual da escola de Direito e \textit{rank} é a classificação da escola de Direito (com \textit{rank = 1} sendo o melhor posto de classificação).
\begin{enumerate}[label=(\alph*)]
    \item Explique a razão por que esperamos $\beta_5 \leq 0$.
    \item Quais são os sinais que você espera para os outros parâmetros de inclinação? Justifique sua resposta.
    \item Utilizando os dados do arquivo LAWSCH85, a equação estimada é:
    \begin{equation*}\begin{aligned}
        \widehat{\log(salary)} =\; & 8,33 + 0,0047LSAT + 0,248colGPA +0,095\log(libvol)\\ & + \; 0,038\log(cost) -0,0033rank
        \end{aligned}
    \end{equation*} \begin{equation*}
        n = 136, \ \ \ \ R^2 = 0,842
    \end{equation*} Qual é a diferença \textit{ceteris paribus} prevista no salário para as escolas com um \textit{colGPA} mediano diferente em um ponto? (Descreva em resposta percentual.)
    \item Interprete o coeficiente da variável $\log(libvol)$.
    \item Você diria que é mais razoável frequentar uma escola de Direito que tem uma classificação melhor? Qual é a diferença no salário inicial esperado para uma escola que tem uma classificação igual a 20?
\end{enumerate}
\begin{solution}
\begin{enumerate}[label=(\alph*)]
    \item Espera-se \(\beta_5\leq 0\) porque quanto pior a classificação da escola (rank maior), menor tende a ser o salário inicial esperado. Reflete uma penalização salarial associada às faculdades de \textit{ranks} mais baixos.
    Se \(\beta_5 = 0\), então o ranking da escola não tem efeito \textit{ceteris paribus} sobre o salário inicial esperado. Isso significa que, controlando pelas demais variáveis do modelo, a posição da escola no ranking não está associada a diferenças salariais.
    \item Espera-se \(\beta_1>0\), \(\beta_2>0\), \(\beta_3>0\) e \(\beta_4>0\), pois maiores valores de \textit{LSAT} e \textit{colGPA} indicam melhor desempenho acadêmico, enquanto maiores valores de \textit{libvol} e \textit{cost} refletem maior qualidade e recursos das instituições. Assim, ceteris paribus, essas características tendem a estar associadas a salários iniciais mais elevados.
    \item Relação log-lin. 24,8\%.
    \item Elasticidade (log-log): um aumento de 1\% em \textit{libvol} está associado a um aumento de 0,095\% no salário inicial esperado.
    \item Sim, mas o efeito marginal de \textit{rank} não é muito grande. Uma piora em 20 posições (+20) reduz \(\log(salary)\) em
    \[
     \Delta x \hat\beta\cdot100 =20(-0{,}0033)\cdot 100=-6,6
    \]
    o que corresponde a uma queda aproximada de 6,6\% no salário inicial esperado.
\end{enumerate}
\end{solution}

\question (3.CS \citep{wooldridge2016}) Use os dados do arquivo \textit{HPRICE1} para estimar o modelo
\[
price=\beta_0+\beta_1sqrft+\beta_2bdrms+u
\]
em que \textit{price} é o preço da residência em milhares de dólares.
\begin{enumerate}[label=(\alph*)]
    \item Escreva os resultados em forma de equação.
    \item Qual é o aumento estimado do preço de uma casa com um quarto a mais, mantendo a área constante?
    \item Qual é o aumento estimado do preço de uma casa com um quarto adicional de 140 pés quadrados? Compare isso com sua resposta do item (b).
    \item Qual porcentagem de variação do preço é explicada pela área e número de quartos?
    \item A primeira casa da amostra tem \(sqrft=2438\) e \(bdrms=4\). Encontre o preço de venda previsto dessa casa a partir da linha de regressão MQO.
    \item O preço de venda real da primeira casa da amostra foi US\$300.000 (assim, \(price=300\)). Encontre o resíduo para esta casa. O valor sugere que o comprador pagou pouco ou muito pela casa?
\end{enumerate}


\end{questions}
\bibliographystyle{apalike}
\bibliography{refs}
\end{document}